\subsection{window-\texorpdfstring{$6$}{6} 的刚性子覆盖、coinvariant 热交叉与 Fibonacci 双标架}\label{subsec:conclusion-window6-subcover-coinvariant-thermal-resonance-lock}

沿用定理 \ref{thm:conclusion-window6-minimal-step-suffix-suspension}、
定理 \ref{thm:conclusion-window6-full-center-derived-central-quotient}、
定理 \ref{thm:conclusion-window6-multiplicity-pinning}、
定理 \ref{thm:conclusion-window6-permutation-module-visible-hidden-splitting}、
定理 \ref{thm:conclusion-window6-hidden-a-type-weyl-package}、
定理 \ref{thm:conclusion-window6-hidden-reflection-invariant-polynomial-ring}、
推论 \ref{cor:conclusion-window6-fiber-gauge-laplacian-spectral-functions}、
定理 \ref{thm:conclusion-window6-visible-crt-arithmetic-phase-space}、
定理 \ref{thm:foldbin-gauge-invariant-ring}、
定理 \ref{thm:foldbin-coinvariant-gorenstein} 与
定理 \ref{thm:fold-critical-resonance-constant} 的记号。前述链条已经分别把
window-\(6\) 的最小二重点层、dyadic 全中心、微态置换模裂分、隐藏反射包、
Laplacian 谱闭式以及可见 CRT 算术相空间固定下来。把这些几何、代数、表示与算术信息并读之后，还可进一步压缩出一条更紧的统一链：最小退化壳不是松散边层，而是由尾部 \(\mathtt{01}\) rigid tail 切出的 \(X_4\)-子覆盖；完整中心并不止于 boundary 的三维可见荷，而是再附带五维的 cyclic 最小壳余量；全部 gauge-invariant 微态多项式既分裂为 visible 线性层与隐藏不变量层，又把所有非平凡剩余压缩进一个顶次数为 \(74\) 的有限 coinvariant 包；与此同时，fold 侧的 Fibonacci 共振双频与 bin-fold 侧的三相阈值在 \(V_6\)-坐标上严格重合，从而把根截面、中心裂分、不变量代数与有限温热交叉统一为同一对 \((8,13)\) 控制的有限对象。

\begin{theorem}[最小二重点层的刚性子覆盖与 \texorpdfstring{$B_3$}{B3} 截面]\label{thm:conclusion-window6-minimal-shell-rigid-subcover-root-slice}
\leanverified{paper\_conclusion\_window6\_minimal\_shell\_rigid\_subcover\_root\_slice}
定义后缀嵌入
\[
\iota_4:X_4\hookrightarrow X_6,
\qquad
\iota_4(x):=x\mathtt{01}.
\]
则有精确恒等式
\[
S_{6,2}=\iota_4(X_4).
\]
进一步，
\[
S_{6,2}\cap X_6^{\mathrm{bdry}}
=
X_6^{\mathrm{bdry}}
=
\{\mathtt{100001},\mathtt{100101},\mathtt{101001}\},
\]
\[
S_{6,2}\cap X_6^{\mathrm{cyc}}
=
\{\mathtt{000001},\mathtt{010001},\mathtt{001001},\mathtt{000101},\mathtt{010101}\}.
\]
在定理 \ref{thm:conclusion-window6-multiplicity-pinning} 的 \(B_3\) pinned root datum 下，
若记
\[
L_2:=\{\alpha\in\Phi_{\mathrm{long}}(B_3):\ \alpha_1=0\},
\]
则上述五个 cyclic 元恰对应
\[
\{-e_1\}\sqcup L_2.
\]
因此
\[
Z(\mathcal G_6^{\mathrm{bin}})
=
Z_{\mathrm{cyc}}\oplus Z_{\partial}
\cong
(\ZZ/2\ZZ)^5\oplus(\ZZ/2\ZZ)^3
\cong
(\ZZ/2\ZZ)^8,
\]
其中 \(Z_{\mathrm{cyc}}\) 由上述五条 cyclic 二点纤维生成，而 \(Z_{\partial}\) 正是 boundary sheet parity 的规范三维中心。
\end{theorem}
\begin{proof}
定理 \ref{thm:conclusion-window6-minimal-step-suffix-suspension} 已经给出
\[
S_{6,2}=\{x\mathtt{01}:\ x\in X_4\},
\]
故第一式成立。

再由定理 \ref{thm:conclusion-window6-full-center-derived-central-quotient}，最小二重点层在 boundary 与 cyclic 两部分的精确分裂为
\[
X_6^{\mathrm{bdry}}
=
\{\mathtt{100001},\mathtt{100101},\mathtt{101001}\},
\]
与
\[
\{\mathtt{000001},\mathtt{010001},\mathtt{001001},\mathtt{000101},\mathtt{010101}\}.
\]
而定理 \ref{thm:conclusion-window6-multiplicity-pinning} 把全部 \(d_6^{\mathrm{bin}}=2\) 的 cyclic 根精确钉扎为
\[
\{-e_1\}\sqcup L_2.
\]
于是最小壳的 cyclic 部分正是一条 codimension-one 的长根截面，再附加唯一短根 \(-e_1\)。

最后，定理 \ref{thm:conclusion-window6-full-center-derived-central-quotient} 已证明
\[
Z(\mathcal G_6^{\mathrm{bin}})
=
Z_{\mathrm{cyc}}\oplus Z_{\partial}
\cong
(\ZZ/2\ZZ)^5\oplus(\ZZ/2\ZZ)^3,
\]
其中 \(Z_{\mathrm{cyc}}\) 由五条 cyclic 二点纤维生成，\(Z_{\partial}\) 由三条 boundary 二点纤维生成。证毕。
\end{proof}

\begin{theorem}[微态 gauge-invariant 代数的自由生成度分布]\label{thm:conclusion-window6-full-gauge-invariant-polynomial-ring}
\leanverified{paper\_conclusion\_window6\_full\_gauge\_invariant\_polynomial\_ring}
记
\[
M_6:=\CC[\Omega_6],
\qquad
P_6:=\operatorname{Sym}(M_6^\ast),
\qquad
G_6^{\mathrm{bin}}:=\mathcal G_6^{\mathrm{bin}}.
\]
则存在代数独立的齐次生成元
\[
\ell_1,\dots,\ell_{21},
\qquad
u_{2,1},\dots,u_{2,21},
\qquad
u_{3,1},\dots,u_{3,13},
\qquad
u_{4,1},\dots,u_{4,9},
\]
满足
\[
\deg \ell_i=1,\qquad
\deg u_{2,j}=2,\qquad
\deg u_{3,k}=3,\qquad
\deg u_{4,r}=4,
\]
并使得
\[
P_6^{G_6^{\mathrm{bin}}}
\cong
\CC[\ell_1,\dots,\ell_{21},u_{2,1},\dots,u_{2,21},u_{3,1},\dots,u_{3,13},u_{4,1},\dots,u_{4,9}]
\]
成为自由多项式环。等价地，
\[
\mathrm{Hilb}\!\left(P_6^{G_6^{\mathrm{bin}}},t\right)
=
\frac{1}{(1-t)^{21}(1-t^2)^{21}(1-t^3)^{13}(1-t^4)^9}.
\]
特别地，代数独立的基本 gauge-invariant 生成元总数恰为
\[
21+21+13+9=64=\dim_{\CC}M_6.
\]
\end{theorem}
\begin{proof}
由定理 \ref{thm:conclusion-window6-permutation-module-visible-hidden-splitting}，
\[
M_6=M_6^{\mathrm{vis}}\oplus M_6^{\mathrm{hid}}
=
M_6^{\mathrm{vis}}\oplus \mathscr R_6,
\]
其中 \(G_6^{\mathrm{bin}}\) 在 \(M_6^{\mathrm{vis}}\) 上平凡作用。故
\[
P_6^{G_6^{\mathrm{bin}}}
\cong
\operatorname{Sym}\!\bigl((M_6^{\mathrm{vis}})^\ast\bigr)\otimes
\operatorname{Sym}(\mathscr R_6^\ast)^{G_6^{\mathrm{bin}}}.
\]

取 \((M_6^{\mathrm{vis}})^\ast\) 的一组基
\[
\ell_1,\dots,\ell_{21},
\]
便有
\[
\operatorname{Sym}\!\bigl((M_6^{\mathrm{vis}})^\ast\bigr)
\cong
\CC[\ell_1,\dots,\ell_{21}].
\]
另一方面，定理 \ref{thm:conclusion-window6-hidden-reflection-invariant-polynomial-ring} 给出
\[
\operatorname{Sym}(\mathscr R_6^\ast)^{G_6^{\mathrm{bin}}}
\cong
\CC[u_{2,1},\dots,u_{2,21},u_{3,1},\dots,u_{3,13},u_{4,1},\dots,u_{4,9}],
\]
其中生成度恰为
\[
2^{\times 21},\qquad 3^{\times 13},\qquad 4^{\times 9}.
\]
将两部分相乘，即得所述自由多项式结构。

Hilbert 级数由各生成元次数直接相乘得到
\[
\mathrm{Hilb}\!\left(P_6^{G_6^{\mathrm{bin}}},t\right)
=
\frac{1}{(1-t)^{21}(1-t^2)^{21}(1-t^3)^{13}(1-t^4)^9}.
\]
最后，生成元总数为
\[
21+(21+13+9)=64,
\]
而 \(\dim_{\CC}M_6=64\) 由定理 \ref{thm:conclusion-window6-permutation-module-visible-hidden-splitting} 已知。证毕。
\end{proof}

\begin{theorem}[隐藏 coinvariant 包的回文 Gorenstein 结构]\label{thm:conclusion-window6-hidden-coinvariant-gorenstein-package}
在定理 \ref{thm:conclusion-window6-full-gauge-invariant-polynomial-ring} 的记号下，定义隐藏 coinvariant 包
\[
\mathcal C_6
:=
\operatorname{Sym}(\mathscr R_6^\ast)\Big/\Big\langle
\bigl(\operatorname{Sym}(\mathscr R_6^\ast)^{G_6^{\mathrm{bin}}}\bigr)_+
\Big\rangle.
\]
则作为分次向量空间，
\[
P_6
\cong
\operatorname{Sym}\!\bigl((M_6^{\mathrm{vis}})^\ast\bigr)
\otimes
\operatorname{Sym}(\mathscr R_6^\ast)^{G_6^{\mathrm{bin}}}
\otimes
\mathcal C_6.
\]
并且 \(\mathcal C_6\) 是 Artinian Gorenstein 代数，满足
\[
\mathrm{Hilb}(\mathcal C_6,t)
=
(1+t)^{21}(1+t+t^2)^{13}(1+t+t^2+t^3)^9,
\]
\[
\dim_{\CC}\mathcal C_6
=
2^{21}3^{13}4^9
=
2^{39}3^{13}
=
|G_6^{\mathrm{bin}}|,
\]
\[
\deg_{\mathrm{top}}(\mathcal C_6)=74.
\]
特别地，
\[
t^{74}\,\mathrm{Hilb}(\mathcal C_6,t^{-1})
=
\mathrm{Hilb}(\mathcal C_6,t).
\]
\end{theorem}
\begin{proof}
由定理 \ref{thm:conclusion-window6-hidden-a-type-weyl-package}，\(\mathscr R_6\) 正是
\[
W(A_1)^{8}\times W(A_2)^{4}\times W(A_3)^{9}
\]
的规范反射表示。对有限实反射群的反射表示，Chevalley 分解给出分次向量空间同构
\[
\operatorname{Sym}(\mathscr R_6^\ast)
\cong
\operatorname{Sym}(\mathscr R_6^\ast)^{G_6^{\mathrm{bin}}}\otimes \mathcal C_6.
\]
再由
\[
P_6
\cong
\operatorname{Sym}\!\bigl((M_6^{\mathrm{vis}})^\ast\bigr)\otimes
\operatorname{Sym}(\mathscr R_6^\ast),
\]
且 \(G_6^{\mathrm{bin}}\) 在可见因子上作用平凡，便得到题述三张量分解。

另一方面，定理 \ref{thm:foldbin-coinvariant-gorenstein} 在 \(m=6\) 的专门化表明：window-\(6\) 的 coinvariant 代数是 Artinian Gorenstein，并具有 Hilbert 级数
\[
(1+t)^{21}(1+t+t^2)^{13}(1+t+t^2+t^3)^9,
\]
总维数
\[
|G_6^{\mathrm{bin}}|
=(2!)^8(3!)^4(4!)^9
=
2^{39}3^{13},
\]
以及顶次数
\[
8\cdot 1+4\cdot 3+9\cdot 6=74.
\]
由于定理 \ref{thm:conclusion-window6-full-gauge-invariant-polynomial-ring} 已把全部一次不变量完全分离到
\(\operatorname{Sym}((M_6^{\mathrm{vis}})^\ast)\) 因子中，故附录中的 full coinvariant 代数与此处的隐藏 coinvariant 包 \(\mathcal C_6\) 规范一致，从而上述公式直接给出 \(\mathcal C_6\) 的 Hilbert 级数、维数与顶次数。

最后，题述 Hilbert 多项式是若干回文因子
\[
1+t,\qquad 1+t+t^2,\qquad 1+t+t^2+t^3
\]
的乘积，且总次数为 \(74\)，故
\[
t^{74}\,\mathrm{Hilb}(\mathcal C_6,t^{-1})
=
\mathrm{Hilb}(\mathcal C_6,t).
\]
证毕。
\end{proof}

\begin{theorem}[纤维规范 Laplacian 的可见--隐藏平衡温度]\label{thm:conclusion-window6-visible-hidden-thermal-crossover}
\leanverified{paper\_conclusion\_window6\_visible\_hidden\_thermal\_crossover}
沿用定理 \ref{thm:conclusion-window6-fiber-gauge-laplacian-spectrum} 的纤维规范 Laplacian
\[
\Delta_6^{\mathrm{fib}}
\in
\operatorname{End}(M_6),
\]
并定义热核配分函数
\[
Z_6(\beta):=\operatorname{Tr}\!\bigl(e^{-\beta \Delta_6^{\mathrm{fib}}}\bigr)
\qquad (\beta\ge 0).
\]
则有精确闭式
\[
Z_6(\beta)=21+8e^{-2\beta}+8e^{-3\beta}+27e^{-4\beta}.
\]
若再记
\[
Z_6^{\mathrm{vis}}(\beta):=21,
\qquad
Z_6^{\mathrm{hid}}(\beta):=8e^{-2\beta}+8e^{-3\beta}+27e^{-4\beta},
\]
则存在唯一的 \(\beta_c>0\) 使得
\[
Z_6^{\mathrm{vis}}(\beta_c)=Z_6^{\mathrm{hid}}(\beta_c).
\]
等价地，令 \(y=e^{-\beta_c}\in(0,1)\)，则 \(y\) 是四次方程
\[
27y^4+8y^3+8y^2-21=0
\]
在区间 \((0,1)\) 内的唯一根。数值上，
\[
\beta_c\approx 0.212489321431856.
\]
并且
\[
Z_6(\beta_c)=42.
\]
\end{theorem}
\begin{proof}
推论 \ref{cor:conclusion-window6-fiber-gauge-laplacian-spectral-functions} 已给出
\[
Z_6(\beta)
=
\operatorname{Tr}\!\bigl(e^{-\beta \Delta_6^{\mathrm{fib}}}\bigr)
=
21+8e^{-2\beta}+8e^{-3\beta}+27e^{-4\beta}.
\]
因此只需研究函数
\[
f(\beta):=Z_6^{\mathrm{hid}}(\beta)-Z_6^{\mathrm{vis}}(\beta)
=
8e^{-2\beta}+8e^{-3\beta}+27e^{-4\beta}-21.
\]
其导数满足
\[
f'(\beta)=-16e^{-2\beta}-24e^{-3\beta}-108e^{-4\beta}<0,
\]
故 \(f\) 在 \((0,\infty)\) 上严格单调递减。另一方面，
\[
f(0)=8+8+27-21=22>0,
\qquad
\lim_{\beta\to\infty}f(\beta)=-21<0.
\]
由介值定理与严格单调性，存在且仅存在一个 \(\beta_c>0\) 使 \(f(\beta_c)=0\)。

令 \(y=e^{-\beta_c}\)，则 \(0<y<1\) 且
\[
8y^2+8y^3+27y^4=21,
\]
即
\[
27y^4+8y^3+8y^2-21=0.
\]
数值求根给出
\[
\beta_c\approx 0.212489321431856.
\]
在该点上，
\[
Z_6(\beta_c)=Z_6^{\mathrm{vis}}(\beta_c)+Z_6^{\mathrm{hid}}(\beta_c)=21+21=42.
\]
证毕。
\end{proof}

\begin{theorem}[fold 共振双频与 bin-fold 阈值的 Fibonacci 锁定]\label{thm:conclusion-window6-fibonacci-resonance-threshold-lock}
\leanverified{paper\_conclusion\_window6\_fibonacci\_resonance\_threshold\_lock}
在定理 \ref{thm:conclusion-window6-visible-crt-arithmetic-phase-space} 的可见算术相空间识别
\[
X_6\cong \ZZ/(21\ZZ)
\]
之下，同一对 Fibonacci 数
\[
\{F_6,F_7\}=\{8,13\}
\]
同时承担两种互不相同而又严格对齐的角色：
\begin{enumerate}
  \item 在 fold 碰撞谱侧，非平凡主共振双频由
  \[
  k=F_6=8,
  \qquad
  k=F_7=13\equiv -8 \pmod{21}
  \]
  承担；
  \item 在 bin-fold 多重度侧，三相阶梯恰在同一对整数处发生转折，即
  \[
  d_6^{\mathrm{bin}}(w)=4
  \iff
  0\le V_6(w)\le 8,
  \]
  \[
  d_6^{\mathrm{bin}}(w)=3
  \iff
  9\le V_6(w)\le 12,
  \]
  \[
  d_6^{\mathrm{bin}}(w)=2
  \iff
  13\le V_6(w)\le 20.
  \]
\end{enumerate}
因此，window-\(6\) 中的 fold 共振频率与 bin-fold 退化相变并非两套偶然相遇的数值现象，而是由同一对 Fibonacci 共轭数严格支配。
\end{theorem}
\begin{proof}
由定理 \ref{thm:fold-critical-resonance-constant}，fold 碰撞谱中承载非衰减主共振的两条频率始终是
\[
F_m,\qquad F_{m+1}.
\]
取 \(m=6\)，即得
\[
F_6=8,\qquad F_7=13.
\]
又因
\[
F_6+F_7=F_8=21,
\]
故
\[
13\equiv -8\pmod{21}.
\]

另一方面，定理 \ref{thm:conclusion-window6-fibonacci-staircase} 已证明
\[
S_{6,4}=V_6^{-1}(\{0,\dots,8\}),
\qquad
S_{6,3}=V_6^{-1}(\{9,\dots,12\}),
\qquad
S_{6,2}=V_6^{-1}(\{13,\dots,20\}),
\]
这正给出 bin-fold 三相分层在 \(8\) 与 \(13\) 处的两条阈值。两侧合并，结论成立。证毕。
\end{proof}

\begin{corollary}[可见模环的 Fibonacci 三带剖面]\label{cor:conclusion-window6-fibonacci-threeband-profile}
\leanverified{paper\_conclusion\_window6\_fibonacci\_threeband\_profile}
在定理 \ref{thm:conclusion-window6-fibonacci-resonance-threshold-lock} 的识别下，\(\ZZ/(21\ZZ)\) 取代表元 \(0,\dots,20\) 时恰分裂为三条相带
\[
[0,8]\ \sqcup\ [9,12]\ \sqcup\ [13,20].
\]
它们的带宽分别为
\[
9,\qquad 4,\qquad 8,
\]
亦即
\[
F_6+1,\qquad F_5-1,\qquad F_6.
\]
特别地，中带宽度恰为 \(F_5-1\)，而左右两带只差一个单位，反映出
window-\(6\) 的 bin-fold 三相剖面是一种由相邻 Fibonacci 数精确锁死的非对称三带结构。
\end{corollary}
\begin{proof}
由定理 \ref{thm:conclusion-window6-fibonacci-resonance-threshold-lock}，三相区域正是
\[
\{0,\dots,8\},\qquad \{9,10,11,12\},\qquad \{13,\dots,20\}.
\]
故其基数依次为
\[
9,\qquad 4,\qquad 8.
\]
而
\[
9=F_6+1,\qquad 4=F_5-1,\qquad 8=F_6.
\]
证毕。
\end{proof}

\begin{theorem}[最小二重点层对完整 dyadic 中心的规范基向量化]\label{thm:conclusion-window6-minshell-dyadic-center-basis}
沿用定理 \ref{thm:conclusion-window6-minimal-shell-rigid-subcover-root-slice} 与定理 \ref{thm:conclusion-window6-full-center-derived-central-quotient} 的记号。对每个
\[
w\in S_{6,2}
\]
记
\[
F_w:=\bigl(\Fold_6^{\mathrm{bin}}\bigr)^{-1}(w),
\]
并令 \(\tau_w\) 为二点纤维 \(F_w\) 上唯一非平凡的换位、在其余纤维上取恒等。则映射
\[
w\longmapsto \tau_w
\]
给出规范的 \(\FF_2\)-基向量对应
\[
\FF_2^{S_{6,2}}
\xrightarrow{\ \sim\ }
Z\!\bigl(\mathcal G_6^{\mathrm{bin}}\bigr).
\]
更精确地，
\[
S_{6,2}\cap X_6^{\mathrm{cyc}}
\longleftrightarrow
Z_{\mathrm{cyc}}\cong (\ZZ/2\ZZ)^5,
\qquad
S_{6,2}\cap X_6^{\mathrm{bdry}}
\longleftrightarrow
Z_{\partial}\cong (\ZZ/2\ZZ)^3.
\]
因此，window-\(6\) 的完整 dyadic 中心全部集中并可基向量化于刚性子覆盖
\[
S_{6,2}=\iota_4(X_4).
\]
\end{theorem}
\begin{proof}
由定理 \ref{thm:conclusion-window6-minimal-shell-rigid-subcover-root-slice}，
\[
S_{6,2}=\iota_4(X_4)
\]
恰由五条 cyclic 二点纤维与三条 boundary 二点纤维组成。对每条二点纤维 \(F_w\)，其置换群同构于 \(S_2\)，故存在唯一非平凡元素 \(\tau_w\)，并且该元素在直积分解
\[
\mathcal G_6^{\mathrm{bin}}\cong S_2^{\times 8}\times S_3^{\times 4}\times S_4^{\times 9}
\]
中恰落在一个 \(S_2\)-因子上，因此属于中心。

另一方面，定理 \ref{thm:conclusion-window6-full-center-derived-central-quotient} 已给出
\[
Z\!\bigl(\mathcal G_6^{\mathrm{bin}}\bigr)
=
Z_{\mathrm{cyc}}\oplus Z_{\partial}
\cong
(\ZZ/2\ZZ)^5\oplus(\ZZ/2\ZZ)^3.
\]
其中 \(Z_{\mathrm{cyc}}\) 由五条 cyclic 二点纤维生成，\(Z_{\partial}\) 由三条 boundary 二点纤维生成。由于
\[
|S_{6,2}|=8=\dim_{\FF_2}Z\!\bigl(\mathcal G_6^{\mathrm{bin}}\bigr),
\]
这些 \(\tau_w\) 恰构成一组规范基。证毕。
\end{proof}

\begin{theorem}[hidden sector 的三重 \texorpdfstring{$43$}{43} 锁定]\label{thm:conclusion-window6-hidden-triple-fortythree-lock}
\leanverified{paper\_conclusion\_window6\_hidden\_triple\_fortythree\_lock}
在 window-\(6\) 上，以下三种先验异质的 hidden 体量严格相等：
\[
\dim_{\CC}\mathscr R_6=43,
\]
\[
\#\{\text{隐藏非线性自由生成元}\}=21+13+9=43,
\]
\[
\dim_{\CC}M_6^{\mathrm{hid}}
=
\operatorname{mult}(2)+\operatorname{mult}(3)+\operatorname{mult}(4)
=8+8+27=43.
\]
因此，隐藏反射包维数、隐藏不变量生成元总数与 hidden 正谱总重数在 window-\(6\) 上发生严格同值锁定。
\end{theorem}
\begin{proof}
第一式正是定理 \ref{thm:conclusion-window6-hidden-a-type-weyl-package} 的结论。第二式由定理 \ref{thm:conclusion-window6-hidden-reflection-invariant-polynomial-ring} 给出的隐藏不变量次数分布
\[
2^{\times 21},\qquad 3^{\times 13},\qquad 4^{\times 9}
\]
直接相加得到。第三式则由推论 \ref{cor:conclusion-window6-fiber-gauge-laplacian-spectral-functions} 中的 hidden 热核闭式
\[
8e^{-2t}+8e^{-3t}+27e^{-4t}
\]
读出，其正谱总重数为 \(8+8+27\)。又由定理 \ref{thm:conclusion-window6-hidden-a-type-weyl-package}，
\[
M_6^{\mathrm{hid}}\cong \mathscr R_6,
\]
故三者同值。证毕。
\end{proof}

\begin{theorem}[hidden coinvariant 包的独立均匀分解与精确矩]\label{thm:conclusion-window6-hidden-coinvariant-uniform-sum}
沿用定理 \ref{thm:conclusion-window6-hidden-coinvariant-gorenstein-package} 的记号，写
\[
\mathrm{Hilb}(\mathcal C_6,t)=\sum_{k=0}^{74}h_k t^k.
\]
定义随机变量 \(\mathsf D_6\) 的分布为
\[
\PP(\mathsf D_6=k):=\frac{h_k}{\dim_{\CC}\mathcal C_6}
\qquad (0\le k\le 74).
\]
则 \(\mathsf D_6\) 可表示为独立和
\[
\mathsf D_6
\overset{d}{=}
\sum_{i=1}^{21}B_i+\sum_{j=1}^{13}T_j+\sum_{\ell=1}^{9}Q_\ell,
\]
其中
\[
B_i\sim \mathrm{Unif}\{0,1\},
\qquad
T_j\sim \mathrm{Unif}\{0,1,2\},
\qquad
Q_\ell\sim \mathrm{Unif}\{0,1,2,3\},
\]
且彼此独立。特别地，
\[
\EE[\mathsf D_6]=37,
\qquad
\operatorname{Var}(\mathsf D_6)=\frac{151}{6},
\]
并且
\[
\PP(\mathsf D_6=k)=\PP(\mathsf D_6=74-k)
\qquad (0\le k\le 74),
\]
故所有奇数阶中心矩全部为零。
\end{theorem}
\begin{proof}
由定理 \ref{thm:conclusion-window6-hidden-coinvariant-gorenstein-package}，
\[
\mathrm{Hilb}(\mathcal C_6,t)
=(1+t)^{21}(1+t+t^2)^{13}(1+t+t^2+t^3)^9.
\]
把
\[
\frac{1+t}{2},
\qquad
\frac{1+t+t^2}{3},
\qquad
\frac{1+t+t^2+t^3}{4}
\]
分别视为 \(\mathrm{Unif}\{0,1\}\)、\(\mathrm{Unif}\{0,1,2\}\)、\(\mathrm{Unif}\{0,1,2,3\}\) 的概率母函数，便得到所述独立和表示。又
\[
\dim_{\CC}\mathcal C_6
=2^{21}3^{13}4^9
\]
恰好是规范化因子，因此该分布与 Hilbert 系数完全一致。

均匀分布的均值与方差为
\[
\EE[B_i]=\frac12,\qquad \operatorname{Var}(B_i)=\frac14,
\]
\[
\EE[T_j]=1,\qquad \operatorname{Var}(T_j)=\frac23,
\]
\[
\EE[Q_\ell]=\frac32,\qquad \operatorname{Var}(Q_\ell)=\frac54.
\]
由独立性，
\[
\EE[\mathsf D_6]
=21\cdot\frac12+13\cdot 1+9\cdot\frac32
=37,
\]
\[
\operatorname{Var}(\mathsf D_6)
=21\cdot\frac14+13\cdot\frac23+9\cdot\frac54
=\frac{151}{6}.
\]

最后，由同一定理中的回文对称
\[
t^{74}\mathrm{Hilb}(\mathcal C_6,t^{-1})=\mathrm{Hilb}(\mathcal C_6,t)
\]
可知 \(h_k=h_{74-k}\)，故分布关于 \(37=74/2\) 对称，所有奇数阶中心矩因而为零。证毕。
\end{proof}

\begin{corollary}[可见--隐藏热交叉的 quartic-dominated 强化]\label{cor:conclusion-window6-thermal-crossover-quartic-dominated}
\leanverified{paper\_conclusion\_window6\_thermal\_crossover\_quartic\_dominated}
沿用定理 \ref{thm:conclusion-window6-visible-hidden-thermal-crossover} 的记号，令
\[
x_c:=e^{-\beta_c}\in(0,1).
\]
则 \(x_c\) 满足严格夹逼
\[
\frac45<x_c<\frac{13}{16},
\qquad\text{亦即}\qquad
\log\frac{16}{13}<\beta_c<\log\frac54.
\]
并且在该热平衡点上，四次壳的单壳贡献已经单独超过可见层的一半：
\[
27e^{-4\beta_c}>\frac{21}{2}.
\]
因此，window-\(6\) 的可见--隐藏热平衡并不是三壳平均相遇，而是一种 quartic-dominated crossover。
\end{corollary}
\begin{proof}
由定理 \ref{thm:conclusion-window6-visible-hidden-thermal-crossover}，
\[
27x_c^4+8x_c^3+8x_c^2-21=0
\]
且 \(x_c\in(0,1)\) 是该方程的唯一实根。直接计算得
\[
27\Bigl(\frac45\Bigr)^4+8\Bigl(\frac45\Bigr)^3+8\Bigl(\frac45\Bigr)^2-21
=-\frac{453}{625}<0,
\]
\[
27\Bigl(\frac{13}{16}\Bigr)^4+8\Bigl(\frac{13}{16}\Bigr)^3+8\Bigl(\frac{13}{16}\Bigr)^2-21
=\frac{22219}{65536}>0.
\]
由于左端在 \((0,1)\) 上严格递增，故
\[
\frac45<x_c<\frac{13}{16}.
\]
取负对数即得
\[
\log\frac{16}{13}<\beta_c<\log\frac54.
\]

再由 \(x_c>\frac45\) 可得
\[
27x_c^4>27\Bigl(\frac45\Bigr)^4=\frac{6912}{625}>\frac{21}{2},
\]
这正是所述 quartic-dominated 不等式。证毕。
\end{proof}

\begin{theorem}[window-\texorpdfstring{$6$}{6} 的 Fibonacci 六标架锁定]\label{thm:conclusion-window6-fibonacci-sixframe-lock}
定义
\[
N_{6,3}^{\mathrm{hid}}:=13,
\qquad
N_{6,2}^{\mathrm{hid}}:=21
\]
分别为 hidden 非线性自由生成元中三次与二次生成元的个数，并记
\[
\delta_{\partial}:=34
\]
为 boundary 二层 uplift 的唯一非零差值。则有严格锁定
\[
\bigl(
|X_6^{\mathrm{bdry}}|,
\ |S_{6,2}\cap X_6^{\mathrm{cyc}}|,
\ |S_{6,2}|,
\ N_{6,3}^{\mathrm{hid}},
\ |X_6|=N_{6,2}^{\mathrm{hid}},
\ \delta_{\partial}
\bigr)
\]
\[
=
(3,5,8,13,21,34)
=(F_4,F_5,F_6,F_7,F_8,F_9).
\]
换言之，window-\(6\) 同时在 boundary/cyclic 最小壳分裂、隐藏自由生成元次数分布、可见稳定层大小以及 boundary uplift 门槛四个不同界面上实现一条连续的 Fibonacci 六标架。
\end{theorem}
\begin{proof}
由定理 \ref{thm:conclusion-window6-minimal-shell-rigid-subcover-root-slice}，
\[
|X_6^{\mathrm{bdry}}|=3,
\qquad
|S_{6,2}\cap X_6^{\mathrm{cyc}}|=5,
\qquad
|S_{6,2}|=8.
\]
由定理 \ref{thm:conclusion-window6-hidden-reflection-invariant-polynomial-ring}，
\[
N_{6,3}^{\mathrm{hid}}=13,
\qquad
N_{6,2}^{\mathrm{hid}}=21.
\]
而定理 \ref{thm:conclusion-window6-visible-crt-arithmetic-phase-space} 给出
\[
|X_6|=21.
\]
最后，由推论 \ref{cor:terminal-foldbin6-boundary-pure-f9-alias}，
\[
\delta_{\partial}=34=F_9.
\]
于是得到整条数列
\[
3,5,8,13,21,34=(F_4,F_5,F_6,F_7,F_8,F_9).
\]
证毕。
\end{proof}
\leanverified{paper\_conclusion\_window6\_fibonacci\_sixframe\_lock}

\begin{conjecture}[偶数窗口的 subcover--coinvariant--thermal 锁定]\label{conj:conclusion-even-window-subcover-coinvariant-thermal-lock}
设 \(m\) 为充分大的偶数窗口，并假定定理 \ref{thm:conclusion-window6-minimal-step-suffix-suspension} 与猜想 \ref{conj:conclusion-even-window-subcover-center-resonance-lock} 的结构在一般 \(m\) 上延续。则预期还应同时成立：
\begin{enumerate}
\item 最小退化壳 \(S_{m,d_{\min}(m)}\) 规范索引全部 dyadic 中心的一组 \(\FF_2\)-基；
\item 隐藏 coinvariant 包 \(\mathcal C_m\) 仍为 Artinian Gorenstein，并且其顶次数等于相应隐藏反射包的总反射数；
\item 可见/隐藏热核存在唯一交叉参数 \(\beta_c(m)\)；
\item 该交叉点的主导壳由 Fibonacci 共振对 \((F_m,F_{m+1})\) 决定，而不是由全部高壳平均产生。
\end{enumerate}
window-\(6\) 正给出这一强化猜想的首个完全闭合样本：最小壳基向量化由定理 \ref{thm:conclusion-window6-minshell-dyadic-center-basis} 实现，hidden 反射包/不变量/热谱的 \(43\)-锁定由定理 \ref{thm:conclusion-window6-hidden-triple-fortythree-lock} 实现，coinvariant 的热浴模型由定理 \ref{thm:conclusion-window6-hidden-coinvariant-uniform-sum} 实现，而热交叉的 quartic-dominated 主导则由推论 \ref{cor:conclusion-window6-thermal-crossover-quartic-dominated} 实现。
\end{conjecture}


\begin{theorem}[window-\texorpdfstring{$6$}{6} 退化谱的残数三区间与终端 Zeckendorf 窗]\label{thm:conclusion-window6-residual-threeinterval-terminal-window}
在定理 \ref{thm:conclusion-window6-visible-crt-arithmetic-phase-space} 的规范识别
\[
X_6\cong \ZZ/(21\ZZ)\cong \{0,1,\dots,20\}
\]
下，bin-fold 纤维多重度满足严格分段公式
\[
d_6^{\mathrm{bin}}(w)=
\begin{cases}
4,& 0\le V_6(w)\le 8,\\[1mm]
3,& 9\le V_6(w)\le 12,\\[1mm]
2,& 13\le V_6(w)\le 20.
\end{cases}
\]
并且最小退化壳恰为一个终端 Zeckendorf 长度窗：
\[
S_{6,2}=V_6^{-1}(\{13,\dots,20\})
\cong
13+\{0,1,\dots,7\}
\subset \ZZ/(21\ZZ).
\]
换言之，window-\(6\) 的最小壳不是任意八点子集，而是 \(X_4\) 的规范 Zeckendorf 范围在商半环中的终端平移像。
\end{theorem}
\begin{proof}
由定理 \ref{thm:conclusion-window6-fibonacci-resonance-threshold-lock}，window-\(6\) 的 bin-fold 三相阈值精确对应
\[
V_6\in \{0,\dots,8\},\qquad V_6\in \{9,10,11,12\},\qquad V_6\in \{13,\dots,20\},
\]
其多重度分别为 \(4,3,2\)。这给出分段公式。

另一方面，定理 \ref{thm:conclusion-window6-minimal-shell-rigid-subcover-root-slice} 已证明
\[
S_{6,2}=\iota_4(X_4)=\{x\mathtt{01}:x\in X_4\}.
\]
由 Zeckendorf 范围双射 \(V_4(X_4)=\{0,\dots,7\}\)，再注意后缀 \(\mathtt{01}\) 对 \(V_6\) 的贡献恰为 \(F_7=13\)，遂得
\[
V_6(S_{6,2})=13+\{0,\dots,7\}=\{13,\dots,20\}.
\]
证毕。
\end{proof}

\begin{theorem}[window-\texorpdfstring{$6$}{6} 最小退化壳的 \texorpdfstring{$(5+3)$}{(5+3)} 裂分]\label{thm:conclusion-window6-minimal-shell-five-three-cleavage}
记
\[
X_{6,\mathrm{cyc},2}:=X_6^{\mathrm{cyc}}\cap S_{6,2}.
\]
则存在严格分裂
\[
S_{6,2}=X_{6,\mathrm{cyc},2}\sqcup X_6^{\mathrm{bdry}},
\qquad
|X_{6,\mathrm{cyc},2}|=5,
\qquad
|X_6^{\mathrm{bdry}}|=3.
\]
并且在 \(B_3\) 的 pinned root datum 下，
\[
X_{6,\mathrm{cyc},2}\cong \{-e_1\}\sqcup\{\pm e_2\pm e_3\},
\]
而
\[
X_6^{\mathrm{bdry}}=\{\texttt{100001},\texttt{100101},\texttt{101001}\}
\]
恰对应三重零权。故最小退化壳并非单一根层，而是“五个 cyclic 极小根点 + 三个 boundary 零权点”的刚性拼接。
\end{theorem}
\begin{proof}
由定理 \ref{thm:conclusion-window6-minimal-shell-rigid-subcover-root-slice}，
\[
S_{6,2}\cap X_6^{\mathrm{bdry}}
=
X_6^{\mathrm{bdry}}
=
\{\texttt{100001},\texttt{100101},\texttt{101001}\},
\]
且
\[
S_{6,2}\cap X_6^{\mathrm{cyc}}
=
\{\texttt{000001},\texttt{010001},\texttt{001001},\texttt{000101},\texttt{010101}\}.
\]
因此
\[
S_{6,2}=X_{6,\mathrm{cyc},2}\sqcup X_6^{\mathrm{bdry}},
\qquad
|X_{6,\mathrm{cyc},2}|=5,\quad |X_6^{\mathrm{bdry}}|=3.
\]
同一定理还给出 cyclic 二重点的根字典恰为
\[
\{-e_1\}\sqcup\{\pm e_2\pm e_3\},
\]
而定理 \ref{thm:window6-21-adjoint-weight-multiset} 已把三条 boundary 词识别为三重零权。证毕。
\end{proof}

\begin{theorem}[window-\texorpdfstring{$6$}{6} 有限体积 hidden fiber 的积分同调平凡性]\label{thm:conclusion-window6-hiddenfiber-integral-homology-triviality}
对任意有限指标集 \(\Lambda\) 与任意稳定型配置
\[
w=(w_x)_{x\in \Lambda}\in X_6^\Lambda,
\]
定义 cubical realization
\[
\mathcal F(w):=\prod_{x\in\Lambda} I_{d_6^{\mathrm{bin}}(w_x)}
\subset ([0,1]^2)^\Lambda,
\]
其中 \(I_2,I_3,I_4\subset \{0,1\}^2\) 为 window-\(6\) hidden fiber 的三种仿射 Boolean 方块标准形。则
\[
\widetilde H_k(\mathcal F(w);\ZZ)=0
\qquad (k\ge 0),
\]
并且
\[
\chi(\mathcal F(w))=1.
\]
亦即，任意有限体积下的 hidden fiber 虽可在局部组合上呈现 \(2,3,4\) 三档结构，但其整体积分同调始终退化为可缩型。
\end{theorem}
\begin{proof}
由定理 \ref{thm:xi-time-part9xab-window6-hidden-fiber-affine-boolean-normal-form}，每个单点 hidden fiber 都同构于 Boolean 方块 \(\{0,1\}^2\) 在乘积序下的序理想
\[
I_2,\qquad I_3,\qquad I_4
\]
之一。故 \(\mathcal F(w)\) 恰是这些序理想的有限直积。

另一方面，推论 \ref{cor:xi-time-part9xab-window6-hidden-fiber-global-contractibility} 已证明：任意有限体积下的上述全局 cubical realization 都可缩。于是其全部约化积分同调群在非负次数上皆为零，Euler 特征等于 \(1\)。证毕。
\end{proof}

\begin{theorem}[window-\texorpdfstring{$6$}{6} 仿射纤维的稀疏算术支撑]\label{thm:conclusion-window6-affine-fiber-arithmetic-support}
把
\[
\Omega_6=\FF_2^6
\]
视为仿射空间，并记 \(F_w=(\Fold_6^{\mathrm{bin}})^{-1}(w)\)。则 \(21\) 条纤维中恰有 \(11\) 条为仿射子空间，并且
\[
F_w\ \text{为仿射子空间}
\iff
V_6(w)\in \{0,4,8\}\cup\{13,14,15,16,17,18,19,20\}.
\]
更精确地：
\begin{enumerate}
  \item \(\{13,\dots,20\}\) 上的八条纤维全为仿射 \(1\)-平坦；
  \item \(\{0,4,8\}\) 上的三条四重点纤维全为仿射 \(2\)-平坦；
  \item \(\{9,10,11,12\}\) 上四条纤维全为“仿射平面缺一点”；
  \item 其余 \(\{1,2,3,5,6,7\}\) 上六条四重点纤维全为仿射 \(3\)-平坦的非线性真子集。
\end{enumerate}
故仿射可实现性在 window-\(6\) 中并非按纤维大小单调扩展，而是压缩为
\[
\{0,4,8\}\cup\{13,\dots,20\}
\]
这一稀疏算术支撑。
\end{theorem}
\begin{proof}
定理 \ref{thm:terminal-foldbin6-fiber-affine-geometry} 已证明：window-\(6\) 中恰有 \(8\) 条纤维为 affine\(_1\)-flat、\(3\) 条纤维为 affine\(_2\)-flat，其余 \(10\) 条均非仿射子空间；并且三条 affine\(_2\)-flat 恰对应稳定词
\[
\texttt{000000},\qquad \texttt{101000},\qquad \texttt{000010}.
\]
由 Zeckendorf 值直接计算，
\[
V_6(\texttt{000000})=0,\qquad
V_6(\texttt{101000})=1+3=4,\qquad
V_6(\texttt{000010})=8.
\]

另一方面，由定理 \ref{thm:conclusion-window6-residual-threeinterval-terminal-window}，全部二重点纤维恰落在
\[
V_6\in\{13,\dots,20\},
\]
故这八条纤维正对应全部 affine\(_1\)-flat。

余下四条三重点纤维由定理 \ref{thm:conclusion-window6-residual-threeinterval-terminal-window} 必位于
\[
V_6\in\{9,10,11,12\},
\]
而定理 \ref{thm:terminal-foldbin6-fiber-affine-geometry} 表明它们全部属于 non-affine-plane-minus-one-point 型。最后，除去 \(\{0,4,8\}\) 与 \(\{13,\dots,20\}\) 之后的六个残数
\[
\{1,2,3,5,6,7\}
\]
对应全部剩余四重点纤维，它们由同一定理全部属于 non-affine-subset-of-affine\(_3\)-flat 型。证毕。
\end{proof}


\begin{conjecture}[偶数窗口的 subcover--center--resonance 四重锁定]\label{conj:conclusion-even-window-subcover-center-resonance-lock}
对任意充分大的偶数窗口 \(m\)，在可见算术相空间识别
\[
X_m\cong \ZZ/(F_{m+2}\ZZ)
\]
之下，预期以下四条将同时成立：
\begin{enumerate}
  \item 最小退化壳满足终端悬挂公式
  \[
  S_{m,d_{\min}(m)}
  =
  \{x\mathtt{01}:\ x\in X_{m-2}\};
  \]
  \item 因而其基数满足
  \[
  |S_{m,d_{\min}(m)}|=F_m=|X_{m-2}|;
  \]
  \item fold 侧的 Fibonacci 共振双频
  \[
  \{F_m,F_{m+1}\}
  \]
  正是 bin-fold 多重度阶梯的两条转折坐标；
  \item 若 cyclic 扇区带有刚性的 pinned root datum \(\Phi_m\)，则最小壳的 cyclic 部分应当切出 \(\Phi_m\) 的一张 codimension-one 根截面，而 dyadic 中心相应分裂为
  \[
  Z(\mathcal G_m^{\mathrm{bin}})
  =
  Z_{m,\partial}\oplus Z_{m,\mathrm{cyc,min}},
  \]
  其中 \(Z_{m,\mathrm{cyc,min}}\) 由最小 cyclic 二点纤维生成。
\end{enumerate}
window-\(6\) 给出该猜想的第一个完全闭合样本：第 \(1\) 条与第 \(2\) 条由定理 \ref{thm:conclusion-window6-minimal-shell-rigid-subcover-root-slice} 验证，第 \(3\) 条由定理 \ref{thm:conclusion-window6-fibonacci-resonance-threshold-lock} 验证，而第 \(4\) 条则由定理 \ref{thm:conclusion-window6-minimal-shell-rigid-subcover-root-slice} 的 \(B_3\) 截面与中心裂分直接实现。特别地，这一猜想把猜想 \ref{conj:conclusion-even-window-minimal-step-terminal-suspension} 的终端悬挂律，推进为一个同时联结最小壳、完整中心、根截面与 Fibonacci 共振对的跨窗口统一原理。
\end{conjecture}

\paragraph{统一读法}
补上最小壳的残数—拓扑—仿射层后，window-\(6\) 可进一步压缩为一个更细的有限统一体：
\[
\text{刚性 }X_4\text{ 子覆盖}
\Longleftrightarrow
\text{残数三区间与终端 Zeckendorf 窗}
\Longleftrightarrow
\text{最小壳的 }(5+3)\text{ 根云/边界裂分}
\Longleftrightarrow
\text{hidden fiber 的 cubical 可缩性与仿射稀疏支撑}
\Longleftrightarrow
\text{隐藏中心余量}
\Longleftrightarrow
\text{自由不变量环}
\Longleftrightarrow
\text{coinvariant 有限包}
\Longleftrightarrow
\text{热平衡交叉}
\Longleftrightarrow
\text{Fibonacci 共振双标架}.
\]
在这一统一体中，几何侧给出 \(B_3\) 根截面、boundary/cyclic 分层以及最小壳的终端窗口化；hidden 侧把每条局部纤维 rigidify 为二维 Boolean 方块的序理想，并在任意有限体积上继续保持可缩；仿射算术侧则把 \(21\) 条纤维压成
\(\{0,4,8\}\cup\{13,\dots,20\}\) 这一稀疏支撑。与此同时，代数侧给出
\((\ZZ/2\ZZ)^8\) 的完整 dyadic 中心与 \(64\) 个自由 gauge-invariant 生成元；表示论侧把全部隐藏非平凡性压缩进一个顶次数为 \(74\) 的 Artinian Gorenstein coinvariant 包；热侧则在完全有限维的 Laplacian 谱上出现唯一的可见--隐藏平衡温度；而算术侧，fold 共振双频 \((8,13)\) 与 bin-fold 的三相阈值严格重合。因而，window-\(6\) 的真正核心对象不再是单独的根字典、单独的纤维直方图或单独的 Fourier 共振，而是这些 rigid 机制在同一个有限窗口上的同步闭合。

\endinput
